\chapter{Complexity — Reductions}
%%% generated by the blueprint pipeline from TCSlib.Complexity.SpaceComplexity.Logspace.Reductions
\section{Overview}

Declarations of the Lean module \texttt{TCSlib.Complexity.SpaceComplexity.Logspace.Reductions}.

\section{Declarations}

\begin{definition}[Logspace many-one reducibility]
\lean{Complexity.LogspaceReducible}
\statementsource{BGS75}{pdf-2459e3840724-p004-b007}
\label{Complexity.LogspaceReducible}
\leanok
\uses{Complexity.ImplicitlyLogspaceComputable}
For Boolean languages $B$ and $C$, say that $B$ is logspace reducible to $C$ if there is
a function $f$ from Boolean lists to Boolean lists that is implicitly computable in
logarithmic space and, for every Boolean list $x$, $x$ belongs to $B$ if and only if
$f(x)$ belongs to $C$.
\end{definition}

\begin{definition}[Nondeterministic logspace completeness]
\lean{Complexity.NLComplete}
\statementsource{Vad12}{pdf-bf2ce5422088-p051-b002}
\label{Complexity.NLComplete}
\leanok
\uses{Complexity.NL}
A Boolean language $C$ is complete for nondeterministic logarithmic space if $C$ belongs
to nondeterministic logarithmic space and every language $B$ in that class is logspace
reducible to $C$.
\end{definition}

\begin{theorem}[Polynomial length bound under composition]
\lean{Complexity.comp_length_le}
\label{Complexity.comp_length_le}
\leanok
\difficulty{4}
Let $f$ and $g$ be functions on bit strings, and let $C,c,D,d$ be natural numbers.
Suppose that for every bit string $x$, $|f(x)|\langle C(|x|+1)^c$, and for every bit string
$y$, $|g(y)|\langle D(|y|+1)^d$. Then for every $x$, $|g(f(x))|\langle D(C+1)^d(|x|+1)^{cd}$.
\end{theorem}

\begin{theorem}[Composition of implicit logspace computations]
\lean{Complexity.ImplicitlyLogspaceComputable.comp}
\label{Complexity.ImplicitlyLogspaceComputable.comp}
\leanok
\uses{Complexity.ImplicitlyLogspaceComputable, Complexity.ImplicitlyLogspaceComputable.compQ, Complexity.compQ, Complexity.comp_length_le, Complexity.indexLang, Complexity.mem_LOGSPACE_preimage, Complexity.mem_indexLang_compQ_iff}
\difficulty{4}
Suppose $f$ and $g$ map Boolean words to Boolean words, each has output length bounded
by a polynomial in the input length, and each of the following two languages is
decidable in logarithmic space: the encoded pairs $\langle
x,\operatorname{bits}(i)\rangle$ for which the $i$-th output bit is true, and the
encoded pairs for which $i$ is below the output length. Then $g\circ f$ also has
polynomially bounded output length and both of its corresponding index languages are
decidable in logarithmic space.
\end{theorem}

\begin{theorem}[Transitivity of Logspace Reducibility]
\lean{Complexity.LogspaceReducible.trans}
\label{Complexity.LogspaceReducible.trans}
\leanok
\uses{Complexity.ImplicitlyLogspaceComputable.comp}
\difficulty{3}
For Boolean languages $B$, $C$, and $D$, if $B$ is logspace reducible to $C$ and $C$ is
logspace reducible to $D$, then $B$ is logspace reducible to $D$.
\end{theorem}

\begin{theorem}[Closure of log space under reductions]
\lean{Complexity.mem_LOGSPACE_of_logspaceReducible}
\label{Complexity.mem_LOGSPACE_of_logspaceReducible}
\leanok
\uses{Complexity.LOGSPACE, Complexity.mem_LOGSPACE_preimage}
\difficulty{3}
Let $B$ and $C$ be languages of finite binary words. If $B$ is log-space reducible to
$C$ and $C$ is decidable in space bounded by a constant multiple of $\log_2(n)+1$, then
$B$ is also decidable in space bounded by a constant multiple of $\log_2(n)+1$.
\end{theorem}

\begin{theorem}[Logspace Reducibility Implies Polynomial-Time Reducibility]
\lean{Complexity.LogspaceReducible.polyTimeReducible}
\label{Complexity.LogspaceReducible.polyTimeReducible}
\leanok
\uses{Complexity.ImplicitlyLogspaceComputable.polyTimeComputable}
\difficulty{1}
For Boolean languages $B$ and $C$, if $B$ is logspace reducible to $C$, then $B$ is
polynomial-time reducible to $C$.
\end{theorem}

\begin{theorem}[A complete language in log space collapses the classes]
\lean{Complexity.NL_eq_LOGSPACE_of_nlComplete_mem_LOGSPACE}
\label{Complexity.NL_eq_LOGSPACE_of_nlComplete_mem_LOGSPACE}
\leanok
\uses{Complexity.LOGSPACE, Complexity.LOGSPACE_subset_NL, Complexity.NL, Complexity.NLComplete, Complexity.mem_LOGSPACE_of_logspaceReducible}
\difficulty{2}
Let $C$ be a language over the Boolean alphabet. If $C$ belongs to $\mathrm{NL}$ and
every language in $\mathrm{NL}$ is log-space reducible to $C$, and if $C$ belongs to
deterministic logarithmic space, then $\mathrm{NL}=\mathrm{LOGSPACE}$.
\end{theorem}
